\documentclass[11pt]{article}

\usepackage[T1]{fontenc}
\usepackage{amsmath,amssymb}
\usepackage{booktabs,tabularx,array}
\usepackage[margin=1in]{geometry}
\usepackage{microtype}
\usepackage{enumitem}
\usepackage{xcolor}
\usepackage{hyperref}
\usepackage{titlesec}

\definecolor{irisblue}{HTML}{22577A}
\definecolor{irisgray}{HTML}{4B5563}
\hypersetup{
  colorlinks=true,
  linkcolor=irisblue,
  urlcolor=irisblue,
  pdftitle={A Complete-F Self-Gated Power-Mean--Anchored Hill-q Rule},
  pdfauthor={IRIS methodological development memo}
}
\titleformat{\section}{\Large\bfseries\color{irisblue}}{\thesection}{0.65em}{}
\titleformat{\subsection}{\large\bfseries\color{irisgray}}{\thesubsection}{0.65em}{}
\setlength{\parindent}{0pt}
\setlength{\parskip}{0.65em}
\setlist{itemsep=0.25em,topsep=0.35em}

\newcommand{\sigmoid}{\operatorname{sigmoid}}
\newcommand{\LSE}{\operatorname{LSE}}

\title{\textbf{A Complete-$F$ Self-Gated Power-Mean--Anchored\\Hill-$q$ Rule for Long-Peptide Aggregation}}
\author{IRIS methodological development memo}
\date{August 24, 2026}

\begin{document}
\maketitle

\begin{abstract}
IRIS assigns Level~1 scores to candidate minimal-epitope--HLA tuples, whereas
the PDAC and SARS-CoV-2 assays label unresolved long peptides. The first
self-gated Hill-$q$ proposal was anchored at the strongest candidate and could
therefore express only scores between maximum and log-sum-exp. New full-roster
tournament output suggested that COVID favors below-maximum averaging, but
those outputs also exposed a more fundamental problem: candidates omitted from
Level~1 computation had been represented by the numerical floor. The apparent
preference therefore partly measures enumeration coverage rather than biology.
This revision makes complete Level~1 tensors a prerequisite, replaces the
maximum anchor by a generalized power-mean anchor, retains floor-aware Hill
corroboration and the self-consistent gate, and makes worst-cohort rank the
primary cross-context calibration objective. No production parameter grid is
frozen here. The complete-$F$ rerun and clean-data diagnosis must occur first.
\end{abstract}

\section{Why Level~2 aggregation is needed}

\subsection{The experimental unit is a long peptide}

For a candidate minimal epitope $e_i$ and restricting HLA allele $h_i$, IRIS
produces a nonnegative mechanistic score
\begin{equation}
F_i=F(e_i,h_i).
\end{equation}
The definition and calibration of $F$ belong to Level~1. The Level~2 problem
arises because a measured long peptide $\ell$ generally contains a roster
\begin{equation}
\mathcal C_\ell=\{(e_i,h_i):i=1,\ldots,K_\ell\}
\end{equation}
of candidate minimal-epitope--HLA tuples, while the observed PDAC or COVID
response label belongs to $\ell$. The experiment does not identify which
member of $\mathcal C_\ell$, if any, caused the response. Level~2 must
therefore map the complete candidate-score roster to one endpoint score:
\begin{equation}
S_\ell=S(F_1,\ldots,F_{K_\ell}).
\end{equation}

Because the Level~1 score is nonnegative and multiplicative, the computation
uses the log scale
\begin{equation}
z_i=\log(F_i+\epsilon),\qquad \epsilon=10^{-12}.
\end{equation}
This transformation stabilizes products and makes score ratios additive.
Neither $F_i$ nor $\sigmoid(z_i)$ is a response probability. The aggregate is
a ranking score on the Level~1 scale, not a calibrated probability of
immunogenicity.

\subsection{Three aggregation questions}

Three plausible biological configurations motivate the design.

\textbf{One-candidate sufficiency.}
One candidate may be strong enough to explain the long-peptide response. A
rule should be able to protect such a leader from dilution by weak roster
members.

\textbf{Leader skepticism.}
An isolated high score may be less persuasive when the rest of the roster is
weak. A rule should also be able to average below the leader when the
component/cohort geometry supports that behavior.

\textbf{Multi-candidate corroboration.}
Several consequential candidates may jointly provide persuasive evidence even
when none is individually decisive. Their magnitudes should be allowed to
accumulate, but not merely because more computational opportunities were
enumerated.

Three familiar fixed rules isolate these questions:
\begin{align}
S_{\max}(z)&=\max_i z_i,\\
S_{\mathrm{mean}}(z)&=\frac{1}{K_\ell}\sum_i z_i,\\
S_{\LSE}(z)&=\log\sum_i e^{z_i}.
\end{align}
Maximum protects a leader but discards all corroboration. Mean discounts an
isolated leader but can be dominated by genuinely weak candidates.
Log-sum-exp accumulates evidence but rises with candidate opportunity. The
revised adaptive family first chooses an anchor on the mean-to-maximum axis,
then offers corroboration toward a complete accumulation ceiling, and finally
lets the developing aggregate gate that offer.

\subsection{Roster geometry and why coverage matters}

Candidate opportunity differs before any score is computed. A peptide of
length $L\geq12$ contains
\begin{equation}
W(L)=\sum_{r=9}^{12}(L-r+1)=4L-38
\end{equation}
sequence windows of lengths 9--12 before HLA expansion. Thus a 15-mer offers
22 sequence windows while a 27-mer offers 70. Pairing those windows with a
patient's distinct HLA alleles expands the tuple roster further. Peptide
length, HLA coverage, homozygosity, overlapping frames, and preprocessing
decisions can therefore change the opportunity set seen by Level~2.

This does not mean that roster size should automatically increase or decrease
the score. It means that the roster convention must be explicit and that
every declared tuple must be treated consistently. Exact duplicate biological
tuples are removed; distinct minimal epitopes or distinct HLA restrictions
remain distinct. The formula does not claim that the remaining contributions
are statistically independent.

\subsection{What the recent outputs do and do not establish}

The newer full-roster tournament produced two useful warnings. First, under
the candidate-conservative view, the max-anchored adaptive systems had no
survivors, and COVID appeared to favor below-maximum averaging. Second, a
large share of the purported full roster had never received a Level~1 tensor
calculation and was assigned $\log\epsilon$ by omission.

The second fact prevents the first observation from being interpreted
directly. Mean and related operators were partly measuring how many windows
the pipeline skipped. The apparent COVID preference, the zero-survivor result,
and any component-specific magnitudes from that run are therefore provisional.
They motivate a complete-$F$ rerun, not a production alpha grid.

One structural conclusion survives this contamination: a family constrained
to $[\max,\LSE]$ cannot express below-maximum behavior on \emph{any} dataset.
Adding a power-mean anchor repairs that expressive limitation without fitting
its parameter values to the contaminated outputs. Whether complete-$F$ COVID
still prefers a low alpha is an empirical question deliberately left open.

\section{Decision and sequencing}

The current omission-floor full-roster results are not valid calibration data
for a new Level~2 rule. A floor can be interpreted biologically only when the
corresponding Level~1 calculation was actually performed. It cannot stand for
an uncomputed value. Work therefore proceeds in this order:

\begin{enumerate}
  \item Freeze the existing Level~1 geometry, decision regime, exposure grid,
  and NCI-selected settings. Do not reselect Level~1.
  \item Compute the Level~1 $F$ tensor for every declared COVID and PDAC
  9--12-mer--HLA tuple.
  \item Deduplicate the roster as a set of
  (endpoint, minimal epitope, normalized HLA) tuples.
  \item Rebuild and audit the transfer package. Every retained mapping row must
  resolve to one computed tensor score.
  \item Re-diagnose component and Level~2 preferences on complete-$F$ data.
  \item Only then freeze the $\alpha$ grid and jointly select
  $(\alpha,q,c,\kappa)$.
  \item Enter the frozen selected systems into the later
  candidate-conservative tournament.
\end{enumerate}

The omission-floor run may be retained as a sensitivity analysis documenting
coverage-by-omission. It must not supply production parameter choices or be
described as if its inputs were complete.

\section{What changed in the model}

\subsection{Why the maximum anchor was insufficient}

Let $z_i$ be the log-scale Level~1 score for candidate $i$. The former rule
started at $m=\max_i z_i$ and admitted a nonnegative corroboration bonus. It
was therefore bounded by
\[
m\leq S\leq \LSE(z).
\]
No choice of Hill order or gate parameters could reproduce mean behavior below
$m$. That was a structural limitation independent of the data-contamination
problem.

The power-mean anchor introduces the missing downward degree of freedom while
retaining the self-gated corroboration mechanism. On the log-score scale,
$\alpha=0$, $\alpha=1$, and $\alpha=\infty$ recover the arithmetic mean,
log-mean-exp, and maximum, respectively.

\subsection{Complete-$F$ input contract}

For a long peptide $\ell$, let
\[
\mathcal C_\ell=\{(e_i,h_i):i=1,\ldots,K_\ell\}
\]
be the set of unique minimal-epitope--HLA tuples after the 9--12-mer length
restriction. Overlapping frames and different HLA restrictions remain distinct
tuples; duplicate observations of the same tuple enter once.

Let $F_i\geq0$ be the computed Level~1 signal and fix
\[
\epsilon=10^{-12}, \qquad z_i=\log(F_i+\epsilon).
\]
A value $z_i=\log\epsilon$ means that a completed Level~1 calculation produced
zero signal. A missing calculation has no $z_i$ and makes the input package
invalid. This distinction is enforced by the implementation.

\section{Power-mean anchor}

The anchor answers the first Level~2 question: how much should the strongest
candidate dominate before corroboration is considered? It is defined on the
positive scale and then reported on the log scale. For $\alpha>0$,
\[
e^{A_\alpha}
=
\left(\frac{1}{K_\ell}\sum_i(e^{z_i})^\alpha\right)^{1/\alpha}.
\]
Taking the logarithm gives the numerically convenient form below.

For finite $\alpha>0$, define
\begin{equation}
A_\alpha(z)
=
\frac{1}{\alpha}
\log\left(\frac{1}{K_\ell}\sum_{i=1}^{K_\ell}e^{\alpha z_i}\right).
\label{eq:power-anchor}
\end{equation}
The continuous endpoint and maximum limit are
\begin{align}
A_0(z)&=\frac{1}{K_\ell}\sum_i z_i,\\
A_\infty(z)&=\max_i z_i.
\end{align}
The family therefore includes
\begin{align}
A_0&=\operatorname{mean}(z),\\
A_1&=\log\left(\frac{1}{K_\ell}\sum_i e^{z_i}\right)
     =\operatorname{logmeanexp}(z),\\
A_\infty&=\max(z).
\end{align}

The order is monotone:
\[
A_0\leq A_\alpha\leq A_\beta\leq A_\infty
\qquad\text{whenever}\qquad 0\leq\alpha\leq\beta.
\]
Increasing alpha therefore gives progressively more authority to the upper
tail. It does not change the candidate roster and it does not inspect an
endpoint label.

\begin{center}
\small
\begin{tabularx}{0.92\textwidth}{@{}p{0.16\textwidth}>{\raggedright\arraybackslash}X@{}}
\toprule
Anchor order & Interpretation on the log-score roster \\
\midrule
$\alpha=0$ & Arithmetic mean of $z$; an isolated leader can be strongly
discounted by the remaining roster. \\
$0<\alpha<1$ & Smooth interpolation between mean and log-mean-exp. \\
$\alpha=1$ & Log-mean-exp; positive-scale arithmetic mean of $e^z$. \\
$1<\alpha<\infty$ & Increasingly upper-tail-sensitive soft maximum. \\
$\alpha=\infty$ & Exact maximum; reproduces the former anchor. \\
\bottomrule
\end{tabularx}
\end{center}

This anchor is not yet the final aggregate. A low alpha permits descent below
the strongest candidate, while the corroboration stage can still move the
score upward when several completed Level~1 signals support that move. The
roles of alpha and $q$ are therefore different: alpha controls the baseline
attitude toward the leader; $q$ controls how the distribution of positive
signal counts as corroborative breadth.

Only nonnegative $\alpha$ is proposed for the first complete-$F$ evaluation.
Negative orders would give increasing leverage to near-zero candidates and
should not enter without separate scientific justification. The production
set of nonnegative orders is intentionally not specified here.

\section{Floor-aware Hill corroboration}

\subsection{Signal weights}

Recover the floor-free Level~1 signal by
\begin{equation}
x_i=\max(e^{z_i}-\epsilon,0).
\end{equation}
When $X=\sum_i x_i>0$, set $p_i=x_i/X$. Candidates whose completed Level~1
signal is exactly zero carry no Hill mass. This prevents a collection of true
numerical floors from manufacturing breadth.

The $p_i$ sum to one and describe how the completed Level~1 signal is
distributed across candidates. They are compositional weights, not response
probabilities and not probabilities that a particular epitope is causal.

\subsection{Hill effective number}

For order $q\geq0$,
\begin{equation}
N_q(p)=
\begin{cases}
\left(\sum_{i:p_i>0}p_i^q\right)^{1/(1-q)}, & q\neq1,\\[3pt]
\exp\left(-\sum_{i:p_i>0}p_i\log p_i\right), & q=1,
\end{cases}
\end{equation}
with
\[
N_0=\#\{i:x_i>0\},\qquad
N_\infty=\frac{1}{\max_i p_i}.
\]
Define
\begin{equation}
D_q=1-\frac{1}{N_q}.
\label{eq:breadth}
\end{equation}
This breadth fraction is zero for a single positive-signal candidate.
For $K_+$ equal positive signals it is $1-1/K_+$ at every order. The orders
differ only when the roster is uneven. Hill numbers are nonincreasing in $q$:
low orders recognize more of the positive roster as breadth, while high
orders require mass to remain close to the leader.

\begin{center}
\small
\begin{tabularx}{0.92\textwidth}{@{}p{0.16\textwidth}>{\raggedright\arraybackslash}X@{}}
\toprule
Hill order & Breadth interpretation \\
\midrule
$q=0$ & Richness-sensitive: every strictly positive completed signal counts,
regardless of its magnitude. \\
$q=1$ & Entropy-weighted: candidates contribute according to information
mass. \\
$q=2$ & Concentration-sensitive:
$D_2=1-\sum_i p_i^2$. \\
$q=\infty$ & Leader-share-sensitive:
$D_\infty=1-\max_i p_i$. \\
\bottomrule
\end{tabularx}
\end{center}

At $q=0$, adding a tiny but positive signal changes richness. That is an
explicit modeling choice, not an invariance claim. Higher orders reduce this
sensitivity, while exact zeros contribute at no order.

\subsection{Available accumulation}

The complete accumulated ceiling is
\begin{equation}
U=\log\left(\epsilon+\sum_i x_i\right).
\label{eq:ceiling}
\end{equation}
For every nonempty roster and $\alpha\geq0$,
\[
A_\alpha\leq \max_i z_i\leq U.
\]
The offer available above the chosen anchor is
\begin{equation}
C_{\alpha,q}
=D_q\,\max(U-A_\alpha,0).
\label{eq:offer}
\end{equation}
Thus $\alpha$ chooses how skeptical the baseline is about an isolated leader,
while $q$ chooses how concentrated the corroborating signal may be.

The construction can be read in three stages:
\begin{enumerate}
  \item $U-A_\alpha$ asks how much complete accumulation is available above
  the chosen anchor.
  \item $D_q$ asks what fraction of that availability is justified by the
  distribution of positive signal.
  \item $C_{\alpha,q}$ is only an \emph{offer}; the gate below decides how
  much enters the final score.
\end{enumerate}

For a singleton positive-signal roster, $N_q=1$, $D_q=0$, and the offer is
zero. For an all-zero completed roster, the computation exits before
normalization and returns the declared floor. These are different cases:
the first has one informative candidate, while the second has a completed
calculation but no positive Level~1 signal.

\section{Self-consistent gate}

For center $c\in\mathbb R$ and scale $\kappa>0$, define the aggregate as the
solution of
\begin{equation}
\boxed{
S
=
\underbrace{A_\alpha}_{\text{power-mean anchor}}
+
\underbrace{C_{\alpha,q}}_{\text{corroboration offer}}\times
\underbrace{\sigmoid\left(\frac{c-S}{\kappa}\right)}_
{\text{fraction admitted}}
}.
\label{eq:self-gate}
\end{equation}
The gate admits more of the offer when the resulting score remains below the
center and less when the result is already strong.

For a provisional aggregate $a$, write
\begin{equation}
g_{c,\kappa}(a)
=
\sigmoid\left(\frac{c-a}{\kappa}\right)
=
\frac{1}{1+\exp[(a-c)/\kappa]}.
\end{equation}
The center $c$ is the aggregate score at which one-half of the offer is
admitted. The positive width $\kappa$ controls how gradual the transition is:
the gate falls from $0.9$ to $0.1$ over
$2\log(9)\kappa\approx4.39\kappa$ score units.

The definition is implicit by design. Evaluating the gate only at the anchor
would ignore the score increase caused by admitted corroboration. Here a large
offer raises $S$, and the increase in $S$ closes the same gate through which
the offer enters. This negative feedback is the self-gating mechanism.

\begin{center}
\small
\begin{tabularx}{0.94\textwidth}{@{}p{0.22\textwidth}>{\raggedright\arraybackslash}X@{}}
\toprule
Quantity & Question answered \\
\midrule
$A_\alpha$ & How much authority should the roster give its strongest
candidate before corroboration? \\
$U-A_\alpha$ & How much complete positive-scale accumulation is available
above that anchor? \\
$D_q$ & How should the concentration of positive signal count as effective
breadth? \\
$C_{\alpha,q}$ & How much corroboration is offered before sufficiency gating? \\
$g_{c,\kappa}(S)$ & What fraction of the offer does the resulting aggregate
still need? \\
\bottomrule
\end{tabularx}
\end{center}

Let
\[
f(S)=S-A_\alpha-C_{\alpha,q}
\sigmoid\left(\frac{c-S}{\kappa}\right).
\]
At the interval endpoints,
\[
f(A_\alpha)\leq0,\qquad
f(A_\alpha+C_{\alpha,q})\geq0.
\]
Moreover,
\[
f'(S)
=1+\frac{C_{\alpha,q}}{\kappa}
\sigmoid\left(\frac{c-S}{\kappa}\right)
\left[1-\sigmoid\left(\frac{c-S}{\kappa}\right)\right]
>0.
\]
Therefore Equation~\eqref{eq:self-gate} has exactly one solution and certified
bisection may be used on
\begin{equation}
A_\alpha\leq S\leq A_\alpha+C_{\alpha,q}\leq U.
\label{eq:bounds}
\end{equation}

If $K_\ell=0$, or if every completed signal is zero, define
\[
S=\log\epsilon,\qquad C_{\alpha,q}=0.
\]
An empty biological roster is allowed; an incomplete tensor roster is not.

\section{How the rule behaves}

\subsection{Certified bounds and monotonicity}

Equation~\eqref{eq:bounds} says more than that the solver is numerically safe.
It locates the scientific behavior of the score:
\[
A_\alpha
\leq S
\leq A_\alpha+C_{\alpha,q}
\leq U.
\]
The aggregate never falls below its declared power-mean anchor and never
exceeds complete positive-scale accumulation. At $\alpha=\infty$, this becomes
the former max-to-accumulation interval. At $\alpha=0$, the interval begins at
the arithmetic mean of the log scores and therefore includes the missing
below-maximum regime.

For fixed anchor and gate parameters, increasing the offer cannot lower the
solution. Implicit differentiation gives
\begin{equation}
\frac{\partial S}{\partial C}
=
\frac{g_{c,\kappa}(S)}
{1-Cg'_{c,\kappa}(S)}
>0.
\end{equation}
This property is one reason to gate at the solved score rather than evaluate
an explicit shortcut at the fully accumulated endpoint. Alpha changes both
the anchor and the available offer, however, so the final score need not be
monotone in alpha for every roster and gate setting.

\subsection{Limiting cases}

\begin{center}
\small
\begin{tabularx}{0.96\textwidth}{@{}p{0.25\textwidth}>{\raggedright\arraybackslash}X@{}}
\toprule
Limit or roster & Consequence \\
\midrule
One positive candidate & $D_q=0$, so $S=A_\alpha=z_1$ for every alpha,
$q$, $c$, and $\kappa$. \\
All completed signals zero & $S=\log\epsilon$ and no Hill distribution is
formed. \\
$\alpha=0$ & The anchor is $\operatorname{mean}(z)$; the family can discount
an isolated leader. \\
$\alpha=1$ & The anchor is $\operatorname{logmeanexp}(z)$. \\
$\alpha=\infty$ & The anchor is maximum; this recovers the expressive range
of the former proposal, with corrected floor semantics. \\
$c\rightarrow-\infty$ & The gate closes and $S\rightarrow A_\alpha$. \\
$c\rightarrow+\infty$ & The gate opens and
$S\rightarrow A_\alpha+C_{\alpha,q}$. \\
$\kappa\rightarrow\infty$ & The gate tends to one-half and
$S\rightarrow A_\alpha+C_{\alpha,q}/2$. \\
\bottomrule
\end{tabularx}
\end{center}

As $\kappa$ approaches zero, the transition becomes nearly hard:
\begin{equation}
S
\approx
\begin{cases}
A_\alpha, & A_\alpha\geq c,\\
A_\alpha+C_{\alpha,q}, & A_\alpha+C_{\alpha,q}\leq c,\\
c, & A_\alpha<c<A_\alpha+C_{\alpha,q}.
\end{cases}
\label{eq:hard-gate}
\end{equation}
A strong anchor therefore receives essentially no bonus; a weak roster whose
full offer remains below $c$ can receive nearly all of it; and an offer that
would cross $c$ is admitted only until the developing aggregate reaches the
transition region. Finite $\kappa$ joins these cases smoothly.

\subsection{Worked score geometries}

The following examples use log scores far above the numerical floor so that
$x_i\approx e^{z_i}$. They illustrate the decomposition before any production
parameters are chosen.

\begin{center}
\small
\begin{tabularx}{0.98\textwidth}{@{}>{\raggedright\arraybackslash}p{0.17\textwidth}>{\raggedright\arraybackslash}p{0.17\textwidth}>{\raggedright\arraybackslash}X@{}}
\toprule
Roster & Key quantities & Interpretation \\
\midrule
Singleton $(-2)$
& $A_\alpha=-2$; $D_q=0$
& The score is exactly $-2$. A singleton is coherent under every setting. \\
\addlinespace
Three-way weak plateau $(-2,-2,-2)$
& $A_\alpha=-2$; $U\approx-0.901$; $D_q=2/3$
& Alpha and $q$ are irrelevant on an equal plateau. The offer is about
$0.732$, so the gate selects a score in $[-2,-1.268]$. \\
\addlinespace
Isolated leader $(0,-5,-5)$
& $A_0=-3.333$; $A_1\approx-1.085$; $A_\infty=0$;
$U\approx0.013$
& Alpha controls whether the leader is discounted. At $q=2$ the two weak
signals provide little breadth; at $q=0$ their positive richness counts much
more strongly. \\
\addlinespace
Three-way strong plateau $(0,0,0)$
& $A_\alpha=0$; $U\approx1.099$; $D_q=2/3$
& The offer is substantial, but a gate centered below zero can admit only a
small fraction because the anchor is already strong. \\
\addlinespace
All-zero completed roster
& $z_i=\log\epsilon$; $x_i=0$
& The score remains exactly $\log\epsilon$. Repeated numerical floors cannot
manufacture corroboration. \\
\bottomrule
\end{tabularx}
\end{center}

The isolated-leader example displays why both order parameters are needed.
Alpha determines the starting point, while $q$ determines whether the weak
positive tail constitutes breadth. Neither parameter alone spans both
questions. The examples also show why complete-$F$ matters: replacing
uncomputed entries by $\log\epsilon$ would change the low-alpha anchors and
could turn a coverage pattern into an apparent biological geometry.

\section{Calibration and evidence}

\subsection{Joint parameter}

The unit of calibration is the complete tuple
\[
(\alpha,q,c,\kappa).
\]
Parameters may be selected separately for each component model and metric
branch, but each selected tuple is then fixed across every cohort in the
policy. Component choice is not hidden inside the Level~2 fit.

\subsection{Freeze one joint grid}

Selection evaluates one declared Cartesian grid
\begin{equation}
\mathcal G
=
\mathcal A\times\mathcal Q\times\mathcal C\times\mathcal K
\end{equation}
for each component--metric group. Alpha is not screened in a preliminary
stage and then held fixed while the remaining parameters are optimized; doing
so would understate the selection search. Every tuple competes in the same
cohort-specific ranking surface.

The complete grid must be frozen before its empirical surfaces are examined.
Adding intermediate alpha, $q$, $c$, or $\kappa$ values after viewing a
rank-based objective changes that objective through grid density. The selector
therefore has no default alpha list: a production invocation must name the
preregistered values explicitly.

Outputs should retain the complete surfaces, the selected tuple, the alpha
and Hill orders represented among near-optimal tuples, and boundary flags for
$c$ and $\kappa$. A selected endpoint order such as zero or infinity is a
scientifically interpretable member of the frozen set, not an instruction to
extend the grid after the fact.

\subsection{Worst-cohort-first objective}

For each cohort, rank every eligible tuple by the policy metric and convert to
fractional rank. For the all-context policy, select by:

\begin{enumerate}
  \item minimum worst-cohort fractional rank;
  \item minimum equal-cohort mean fractional rank;
  \item minimum equal-cohort mean metric regret;
  \item declared evidence and numerical tie-breaks.
\end{enumerate}

This maximin ordering matches the fact that the candidate-conservative
criterion is defeated by a weak cohort, not rescued by an excellent average.
For the PDAC-only policy the first two quantities reduce to PDAC rank; COVID
remains a frozen transport assessment.

For PR, a tuple remains eligible only when the staged paired-CNAP comparison
supports it directionally over the same component model's fixed maximum in
every policy cohort. ROC uses the same worst-cohort-first ordering on AUROC.
These are parameter-selection rules, not a substitute for the later complete
candidate-conservative tournament.

\subsection{Transport gate}

Leave-one-cohort-out results must be generated before production selection.
The earlier memo treated them only descriptively. A quantitative transport
failure threshold should be preregistered after the complete-$F$ diagnostic
and before the production grid is evaluated. This memo does not invent that
threshold from contaminated outputs.

\subsection{Label blindness and evidential status}

Once a tuple is selected, endpoint adaptation is label blind. The endpoint's
candidate scores determine $A_\alpha$, $D_q$, $C_{\alpha,q}$, and the solved
score. Neither the response label nor the cohort identity chooses a different
alpha, Hill order, or gate for that endpoint.

Labels do participate in model-development selection of
$(\alpha,q,c,\kappa)$. If the same PDAC and COVID cohorts later supply the
directed round robin, that tournament is conditional post-selection evidence,
not independent external confirmation. The selection grid, complete surfaces,
eligibility outcomes, solver contract, input checksums, and this evidential
status belong in the manifest.

\subsection{Component-level scope}

Level~2 selection occurs within each component construction. It can change how
that component's candidate scores are pooled, but it cannot transform a
Full-HLA score geometry into a Mono-Q/full-PN geometry or vice versa. A
component-level context reversal may therefore remain after the best Level~2
rule is chosen. The production report should distinguish failure of the
aggregator to transfer from failure of the underlying component family to
transfer.

\section{Design requirements}

The equations are compact, but they answer only part of the design problem.
The following requirements make the roster, numerical, and evidential
assumptions explicit.

\begin{description}[style=nextline,leftmargin=0pt]
  \item[P1. Complete computation.]
  Every declared roster tuple has a Level~1 tensor result. Missing values fail
  closed.

  \item[P2. Set semantics.]
  The roster is a set of endpoint--minimal-epitope--HLA tuples. Exact
  duplicates do not create breadth.

  \item[P3. Downward and upward reach.]
  The family contains mean, log-mean-exp, and maximum anchors and may admit
  corroboration up to the complete accumulation ceiling.

  \item[P4. One-candidate safety.]
  A singleton positive-signal roster has $D_q=0$, so $S=A_\alpha=z_1$.

  \item[P5. Zero-floor safety.]
  Completed zero signals carry no Hill mass. An all-zero roster cannot
  manufacture corroboration.

  \item[P6. No implicit multiplicity claim.]
  Distinct overlapping windows and HLA tuples may contribute, but the score
  does not assert their biological independence.

  \item[P7. Certified numerics.]
  The unique root is solved only within Equation~\eqref{eq:bounds};
  nonconvergence fails closed.

  \item[P8. Robust calibration.]
  Cross-context selection is worst-cohort-first, with the full tuple selected
  jointly and fixed after selection.

  \item[P9. Honest scope.]
  Level~2 adaptivity cannot repair a component family that is structurally
  wrong for a cohort. Component-level reversals are reported explicitly.
\end{description}

\subsection{Candidate-opportunity comparability}

Roster richness is neither automatically supportive nor automatically
suspect. It may arise from peptide length, HLA coverage, genuine biological
breadth, or computational enumeration. The revised rule exposes rather than
hides its sensitivity:

\begin{itemize}
  \item low alpha anchors respond to the complete roster, including weak
  completed candidates;
  \item $q=0$ treats every positive signal as richness;
  \item higher $q$ values require more signal mass near the leader;
  \item the accumulation gap and self-gate limit how far corroboration can
  raise the anchor;
  \item exact zero signals contribute no Hill mass.
\end{itemize}

Because the low-alpha anchor is intentionally roster-sensitive, complete
enumeration is not merely a data-cleaning preference; it is an identification
condition for the model. Two endpoints cannot be compared fairly if one
received complete Level~1 calculations and the other had uncomputed candidates
inserted at the floor.

The rule does not add a separate penalty based only on $K_\ell$. Such a
penalty would require a stochastic model specifying how extra candidate
opportunities create spurious scores. Without that model it could erase
genuine breadth. Instead, opportunity is handled through a complete declared
roster, explicit alpha and $q$ choices, and sensitivity reporting.

\subsection{Set rather than multiset}

Exact duplicate control occurs before the score vector is formed. A repeated
occurrence of the same endpoint--minimal-epitope--normalized-HLA tuple enters
once, including recurrence caused by HLA homozygosity or duplicated source
rows. The same minimal epitope paired with two distinct HLA alleles forms two
tuples. Overlapping but nonidentical epitopes also remain distinct.

This convention prevents literal duplication from manufacturing breadth. It
does not assert independence among the distinct tuples. Correlated candidates
can still move the aggregate together, and their biological dependence is an
interpretive qualification rather than something the Hill number estimates.

\section{Interpretive qualifications}

\textbf{Anchor order.}
Alpha controls the generalized power-mean baseline. It is not a posterior
probability that the best candidate is causal. A component--metric-specific
alpha remains fixed across policy cohorts after selection.

\textbf{Hill order.}
The order $q$ translates a normalized positive-signal distribution into
effective breadth. It is not a probability of candidate independence and does
not estimate the number of truly causal epitopes.

\textbf{Gate parameters.}
The center $c$ is the aggregate log-$F$ score at which the admitted fraction is
one-half; it is not a response-probability threshold. Its meaning depends on
the frozen definition of $F$, logarithmic transformation, and epsilon. The
width $\kappa$ is a score-scale transition width, not a biological variance.

\textbf{Numerical floor.}
$\log(10^{-12})$ represents a completed Level~1 zero signal or the declared
empty-roster output. It never represents an uncomputed observation in the
replacement workflow.

\textbf{Numerical solution.}
Equation~\eqref{eq:self-gate} is solved as a scalar root on the certified
interval in Equation~\eqref{eq:bounds}. Bisection tolerance, maximum
iterations, and failure behavior are part of the scoring contract. Naive
fixed-point iteration is not part of the definition.

\textbf{Context and causation.}
The aggregate does not use cohort identity or peptide length as inputs. That
label blindness does not prove that its rankings are invariant to assay
geometry, HLA distribution, disease context, or label construction. Transfer
must be tested rather than inferred from the form of the equation.

\textbf{Scope of adaptation.}
The self-gate limits an offer according to the score it produces. It does not
establish that the roster is biologically exhaustive, that distinct tuples are
independent, or that a selected parameter tuple will transfer to an
unrepresented context.

\section{What is and is not frozen}

This revision freezes:
\begin{itemize}
  \item complete-$F$ rather than omission-floor inputs;
  \item exact tuple-set deduplication;
  \item the power-mean anchor in Equation~\eqref{eq:power-anchor};
  \item the floor-aware Hill offer in Equations~\eqref{eq:breadth}--\eqref{eq:offer};
  \item the self-gated equation and certified solver;
  \item joint $(\alpha,q,c,\kappa)$ evaluation;
  \item worst-cohort-first cross-context selection.
\end{itemize}

It deliberately does not freeze:
\begin{itemize}
  \item the production $\alpha$ values;
  \item any conclusion that COVID genuinely prefers mean after complete-$F$
  rescoring;
  \item the production transport-gate threshold;
  \item a selected tuple for any component or metric;
  \item the claim that Level~2 adaptation alone will produce a
  candidate-conservative survivor.
\end{itemize}

\section{Implementation correspondence}

The authoritative implementation is colocated with this memo in
\path{supported_ap_code/crates/select_adaptive_hillq}. It has one
complete-$F$ code path and no compatibility mode for the former max-anchored
or omission-floor selector. Its manifests record the explicit alpha and Hill
grids, the complete formula, the worst-cohort objective, input hashes, solver
settings, and selected parameters.

\section{Post-testing updates and extended formulation}
\label{sec:post-testing-extension}

\subsection{Status and purpose of this update}

The preceding sections record the formulation proposed before complete-$F$
testing. This section records the conclusions reached after that testing. It
does not convert the resulting evidence into independent validation. Rather,
it identifies which structural hypotheses survived, which failed, and how the
retained self-gating mechanism was extended.

The complete-$F$ results did not reject adaptive corroboration. They rejected
the assumption that every short-peptide--HLA tuple should enter one
undifferentiated breadth roster. The working extension retains the maximum as
a sufficiency anchor, measures continuous breadth across distinct short
peptides, and treats evidence through a second distinct HLA as a separate,
bounded source of corroboration.

\subsection{What complete-$F$ testing retained and rejected}

The post-test findings can be summarized as follows:
\begin{enumerate}
  \item Complete Level~1 computation and floor-aware signal semantics remain
  necessary. An uncomputed tensor entry may not be represented by the
  numerical floor.
  \item No common finite-$\alpha$ rule from the original family transported
  across PDAC, SARS-CoV-2 SPIKE, and SARS-CoV-2 NONSPIKE. With the complete
  roster, finite-$\alpha$ anchors were vulnerable to dilution.
  \item Maximum remained a transportable sufficiency anchor. Useful
  corroboration was present at $\alpha=\infty$, although the best absolute
  gate center differed across the three datasets.
  \item A cohort-relative gate was numerically successful but scientifically
  inadmissible because changing the evaluation batch could change the score of
  an endpoint already being scored.
  \item Collapsing the roster to one maximum per distinct short-peptide
  sequence nearly supplied a common rule. A small, separately represented cue
  from the second-best distinct HLA supplied the remaining information.
\end{enumerate}

These findings narrow the three aggregation questions introduced earlier.
One-candidate sufficiency survives through the maximum anchor. The proposed
portable role for leader skepticism did not survive this complete-$F$ test,
so the working rule fixes $\alpha=\infty$. Multi-candidate corroboration
survives, but it must be decomposed into distinct-epitope breadth and
distinct-HLA corroboration.

\subsection{Two projections of the complete candidate roster}

The complete tuple set remains the Level~1 input contract. For one long
peptide, write its aligned scores as $z_{eh}$, where $e$ denotes a distinct
short-peptide sequence and $h$ denotes a normalized HLA identity. Level~2
constructs two endpoint-local projections of this same complete roster.

The distinct-epitope projection is
\begin{equation}
u_e=\max_h z_{eh},
\qquad
\mathcal U=(u_e:e\in\mathcal E),
\label{eq:posttest-epitope-projection}
\end{equation}
where $\mathcal E$ is the set of distinct short-peptide sequences. The
distinct-HLA projection is
\begin{equation}
M_h=\max_e z_{eh},
\qquad
\mathcal M=(M_h:h\in\mathcal H),
\label{eq:posttest-hla-projection}
\end{equation}
where $\mathcal H$ is the endpoint's set of distinct normalized HLAs. In both
projections,
\begin{equation}
m=\max_e u_e=\max_h M_h=\max_{e,h}z_{eh}.
\label{eq:posttest-common-anchor}
\end{equation}

This is projection-aware set semantics. The complete peptide--HLA tensor is
still required, and exact duplicate biological tuples are still removed.
The projections change only the units in which two different kinds of
corroboration are measured. They do not assert statistical independence among
peptides, among HLAs, or between the two projections.

\subsection{Self-consistent distinct-epitope branch}

Apply the floor-aware Hill construction to $\mathcal U$. Define
\begin{align}
x_e&=\max(e^{u_e}-\epsilon,0),
&X_E&=\sum_e x_e,
&p_e&=x_e/X_E,\\
D_E&=1-\max_e p_e,
&U_E&=\log(\epsilon+X_E),
&C_E&=D_E\max(U_E-m,0).
\end{align}
As before, an all-zero roster exits before normalization. The working
distinct-epitope branch score $S_E$ is the unique solution of
\begin{equation}
\boxed{
S_E=m+C_E\sigmoid\left(\frac{c-S_E}{\kappa}\right),
\qquad c=-2.2,\quad \kappa=0.13.
}
\label{eq:posttest-epitope-self-gate}
\end{equation}

This branch is an exact member of the original self-consistent family, with
$\alpha=\infty$ and $q=\infty$, applied to the distinct-epitope projection.
Its negative-feedback interpretation, uniqueness proof, certified bisection
solution, and bounds are inherited unchanged:
\begin{equation}
m\leq S_E\leq m+C_E\leq U_E.
\label{eq:posttest-epitope-bounds}
\end{equation}

\subsection{Second-HLA evidence-qualification branch}

Let $m_2$ be the second-largest member of $\mathcal M$. Define
\begin{equation}
g_H(m_2)=\sigmoid\left(\frac{m_2-t}{\delta}\right),
\qquad
S_H=m+B g_H(m_2).
\label{eq:posttest-hla-gate}
\end{equation}
For the representative working tuple,
\begin{equation}
t=-6.45,
\qquad \delta=0.02,
\qquad B=1.
\end{equation}

The two sigmoids answer different questions. The distinct-epitope branch
controls a roster-dependent accumulation offer $C_E$. Because admitting that
offer changes the score used to judge sufficiency, the gate is evaluated at
the developing score $S_E$: as corroboration raises $S_E$, it closes the same
gate through which it enters. The HLA branch instead asks whether the fixed
evidence $m_2$ qualifies as a credible second presentation route. Adding its
bounded bonus changes neither $m_2$ nor its credibility, so no
self-consistency loop is needed. Making this gate depend on $S_H$ would ask
whether the aggregate still needs HLA corroboration, rather than whether the
second-HLA evidence is credible. Neither sigmoid is a response probability.

\subsection{Extended common-anchor formulation}

The final score is the convex combination
\begin{equation}
S_{\mathrm{hybrid}}=(1-w)S_E+wS_H,
\qquad w=0.12.
\label{eq:posttest-hybrid-convex}
\end{equation}
Using Equation~\eqref{eq:posttest-epitope-self-gate}, this can be written in
anchor--offer--gate language as
\begin{equation}
\boxed{
S_{\mathrm{hybrid}}
=m
+(1-w)\underbrace{C_E
  \sigmoid\left(\frac{c-S_E}{\kappa}\right)}_
  {\text{admitted distinct-epitope corroboration}}
+w\underbrace{B
  \sigmoid\left(\frac{m_2-t}{\delta}\right)}_
  {\text{admitted second-HLA corroboration}}.
}
\label{eq:posttest-hybrid-expanded}
\end{equation}

The maximum $m$ is the common sufficiency anchor. The quantity $C_E$ is the
distinct-epitope corroboration offer, and $B$ is the cap on the second-HLA
offer. The two gates determine the fractions admitted for different reasons,
while $w$ allocates authority between the two corroboration mechanisms. More
precisely, the epitope-breadth uplift is attenuated by $1-w$, and the HLA
mechanism can add at most $wB=0.12$ log-score units. This bound limits score
movement; it does not guarantee that the HLA term cannot determine rankings
among otherwise close endpoints.

\subsection{A pedagogical interpretation and possible higher-HLA extension}

Both branches begin with the strongest predicted short-peptide--HLA route
$m$. They differ only in the biological direction from which they seek
corroboration. In the simplest terms, they ask
\begin{equation}
\boxed{
S_E:\quad \text{Are several distinct short peptides plausible?}
}
\end{equation}
and
\begin{equation}
\boxed{
S_H:\quad \text{Is there a second plausible HLA presentation route?}
}
\end{equation}
The $S_E$ branch adds self-gated corroboration when Level~1 signal is
distributed across multiple distinct short-peptide sequences. The $S_H$ branch
adds bounded corroboration when a second distinct HLA has at least one
sufficiently strong short-peptide candidate. Thus $S_E$ is the strongest route
corroborated along the peptide axis, while $S_H$ is the same strongest route
corroborated along the HLA axis.

The two projections make this symmetry explicit:
\begin{equation}
u_e=\max_h z_{eh}
\qquad\text{versus}\qquad
M_h=\max_e z_{eh}.
\end{equation}
The peptide projection asks how broadly signal is distributed after each
distinct peptide is represented by its strongest HLA. The HLA projection asks
which distinct presentation routes remain credible after each HLA is
represented by its strongest peptide. The second-HLA route may involve either
the peptide attaining the overall maximum or a different peptide. It is not
``the HLA after the one associated with $S_E$,'' because $S_E$ does not have one
unique HLA identity.

This also explains an asymmetry in the possible extensions. The $S_E$ branch
already uses the entire distinct-peptide roster and therefore represents two,
three, or many plausible peptides without separate second- or third-peptide
branches. The present $S_H$ branch uses only the second-largest distinct-HLA
maximum. Higher HLA order statistics therefore provide a natural conceptual
extension. Write
\begin{equation}
M_{(1)}\geq M_{(2)}\geq\cdots\geq M_{(R)}
\end{equation}
for the ordered maxima from distinct normalized HLAs, and define, for
$r\geq2$,
\begin{equation}
S_{H,r}
=m+B_r\sigmoid\left(\frac{M_{(r)}-t_r}{\delta_r}\right).
\label{eq:posttest-higher-hla-branch}
\end{equation}
Here $S_{H,2}$ asks whether a second HLA route is plausible, $S_{H,3}$ asks
whether a third is plausible, and so on. A general multi-HLA extension could take the
form
\begin{equation}
S_R=w_E S_E+\sum_{r=2}^{R}w_rS_{H,r},
\qquad
w_E+\sum_{r=2}^{R}w_r=1,
\label{eq:posttest-multihla-convex}
\end{equation}
or, equivalently,
\begin{equation}
S_R
=m+w_E(S_E-m)
+\sum_{r=2}^{R}w_rB_r
\sigmoid\left(\frac{M_{(r)}-t_r}{\delta_r}\right).
\label{eq:posttest-multihla-expanded}
\end{equation}
The present working rule is recovered with
\begin{equation}
R=2,
\qquad w_E=0.88,
\qquad w_2=0.12.
\end{equation}

This higher-HLA construction is a natural extension of the representation,
not an empirical conclusion from the present tests. Evidence through a third
HLA may be redundant with evidence through the second because HLA-specific
predictions can be correlated. Any future extension should therefore be
frozen before evaluation, define an unavailable higher-order HLA as zero
incremental corroboration, and keep later contributions diminishing and the
total HLA uplift bounded. Possible structural constraints are
\begin{equation}
w_{r+1}B_{r+1}\leq w_rB_r
\end{equation}
and
\begin{equation}
\sum_{r=2}^{R}w_rB_r\leq B_{\mathrm{HLA,max}}.
\end{equation}
These constraints express disciplined hypotheses, not parameter values
supported by the current evidence. Peptide multiplicity is already summarized
continuously by $S_E$; HLA multiplicity is currently truncated at the second
route and therefore admits a natural, but as yet untested, higher-order
extension.

\subsection{Relation to the original self-consistent gate}

The two-branch form is the preferred biological interpretation, but the final
score also has an exact single-gate representation. Let
\begin{equation}
b=B g_H(m_2).
\end{equation}
Then Equation~\eqref{eq:posttest-hybrid-convex} is equivalent to
\begin{equation}
S_{\mathrm{hybrid}}
=A_{\mathrm{eff}}
+C_{\mathrm{eff}}
\sigmoid\left(
  \frac{c_{\mathrm{eff}}-S_{\mathrm{hybrid}}}
       {\kappa_{\mathrm{eff}}}
\right),
\label{eq:posttest-effective-self-gate}
\end{equation}
where
\begin{align}
A_{\mathrm{eff}}&=m+wb,\\
C_{\mathrm{eff}}&=(1-w)C_E,\\
c_{\mathrm{eff}}&=(1-w)c+w(m+b),\\
\kappa_{\mathrm{eff}}&=(1-w)\kappa.
\end{align}
Because $0<w<1$ and $\kappa>0$, this transformed equation has the same
strictly increasing residual and the same unique-root argument as the
original self-consistent gate. The representation is useful for mathematical
correspondence, but it is not preferred for interpretation: its effective
center depends on the endpoint-local HLA evidence and hides the distinction
between the two corroboration sources.

\subsection{Inherited and revised guarantees}

\begin{center}
\small
\begin{tabularx}{0.96\textwidth}{@{}p{0.31\textwidth}>{\raggedright\arraybackslash}X@{}}
\toprule
Property & Post-test status \\
\midrule
Complete-$F$ requirement
& Retained. Every declared peptide--HLA tuple must have a completed Level~1
calculation. \\
Endpoint locality and label blindness
& Retained. Both projections and both gates use only the endpoint's roster
and globally frozen constants. \\
Unique self-consistent solution
& Retained for $S_E$ and inherited by the exact effective representation of the
final score. \\
Maximum lower bound
& Retained: $S_{\mathrm{hybrid}}\geq m$. \\
Complete-accumulation upper bound
& Revised. The epitope branch is bounded by $U_E$, but the direct HLA bonus
means the final hybrid need not be. \\
One-candidate safety
& Refined. One distinct peptide has no epitope-breadth offer, but it may
receive a bounded bonus if it has credible support through a second distinct
HLA. \\
Set semantics
& Refined into distinct-peptide and distinct-HLA projections of the complete
tuple set. \\
Joint frozen calibration
& Expanded to include the projection rules and the HLA-branch parameters. \\
\bottomrule
\end{tabularx}
\end{center}

The direct certified bound for the final score is
\begin{equation}
m\leq S_{\mathrm{hybrid}}
\leq m+(1-w)C_E+wB.
\label{eq:posttest-hybrid-bounds}
\end{equation}
The upper bound is deliberately not identified with $U_E$. The bounded direct
HLA offer was retained because a formulation restricted to the remaining
distance below $U_E$ did not provide a common supported rule.

The production contract must also define edge cases explicitly. An all-zero
completed roster should short-circuit to $\log\epsilon$ before either gate is
formed. If fewer than two distinct normalized HLAs are present, the HLA
corroboration term should be defined by a frozen fallback rather than by an
implicit implementation accident.

\subsection{Post-selection evidence and the next frozen test}

The representative working tuple is
\begin{equation}
(\alpha,q,c,\kappa,t,\delta,B,w)
=(\infty,\infty,-2.2,0.13,-6.45,0.02,1,0.12),
\label{eq:posttest-representative-tuple}
\end{equation}
with the projection rules in
Equations~\eqref{eq:posttest-epitope-projection} and
\eqref{eq:posttest-hla-projection} included as part of the model definition.
Against the same Full-HLA fixed-maximum reference, the staged paired-CNAP
assessment with 600 replications produced:

\begin{center}
\small
\begin{tabular}{@{}lrrr@{}}
\toprule
Dataset & Observed retained difference & Supported magnitude & Survival \\
\midrule
PDAC      & $+0.0003428$ & $+0.0001200$  & $0.9344$ \\
SPIKE     & $+0.0000340$ & $+0.00000589$ & $0.9280$ \\
NONSPIKE  & $+0.0007612$ & $+0.0001587$  & $0.9933$ \\
\bottomrule
\end{tabular}
\end{center}

A local $3^4$ neighborhood varied
\begin{align*}
t&\in\{-6.475,-6.45,-6.425\},
&\delta&\in\{0.015,0.020,0.025\},\\
w&\in\{0.08,0.10,0.12\},
&B&\in\{0.50,0.75,1.00\}.
\end{align*}
Sixty-one of the 81 tuples passed in all three datasets, with passing tuples
at every tested value of every parameter. This demonstrates local numerical
robustness, not 61 independent confirmations. The functional form and its
constants were developed after inspecting the complete-$F$ results, so these
assessments remain post-selection evidence. The next scientific test should
freeze Equation~\eqref{eq:posttest-representative-tuple} before evaluating it
on genuinely held-out data or in any new tournament.

The detailed development record, including unsuccessful endpoint-local forms,
is given in
\path{supported_ap_code/crates/select_adaptive_hillq/full_hla_endpoint_local_hybrid_memo.tex}.
The representative and neighborhood outputs are stored in
\path{artifacts/iris_fullroster_complete_f_20260904/adaptive_rule_exploration/}.

\subsection{Updated implementation contract}

The exploratory implementation is
\path{supported_ap_code/crates/select_adaptive_hillq/src/bin/explore_full_hla_hybrid_breadth.rs}.
It establishes correspondence with the post-test numerical results but has not
replaced the production selector contract. Before the rule is admitted to a
new evaluation, its production implementation should:
\begin{enumerate}
  \item expose a pure endpoint-scoring function with no cohort or batch input;
  \item canonicalize short-peptide sequences and HLA identities before either
  projection is constructed;
  \item impose deterministic ordering before floating-point aggregation;
  \item define empty, all-zero, and fewer-than-two-HLA behavior explicitly;
  \item freeze the Level~1 score definition, logarithmic transformation,
  numerical floor, gate constants, solver tolerance, and failure behavior;
  \item include golden tests that connect the published equations to exact
  implementation outputs; and
  \item record source hashes, input hashes, constants, and model version in
  the evaluation manifest.
\end{enumerate}

The extended rule should therefore be understood as a common-anchor,
dual-corroboration aggregate: a self-consistent distinct-epitope breadth gate
combined with a bounded second-HLA evidence-qualification gate. It preserves
the anchor--offer--gate logic of the original proposal while assigning the two
kinds of breadth to biologically legible units.

\end{document}
